% =============================================================
% Civilytics — LaTeX preamble for Quarto PDF
% Usage in YAML:
%   format:
%     pdf:
%       include-in-header: quarto/latex/civilytics.tex
%       include-before-body: quarto/latex/civilytics-title.tex
% Requires: xelatex or lualatex (for system fonts).
% =============================================================

\usepackage{xcolor}
\usepackage{fontspec}
\usepackage{titlesec}
\usepackage{titling}
\usepackage{fancyhdr}
\usepackage{enumitem}
\usepackage{booktabs}
\usepackage{caption}
\usepackage[hidelinks]{hyperref}
\usepackage{microtype}
\usepackage{tcolorbox}
\tcbuselibrary{skins,breakable}
% xurl must load AFTER hyperref. Without it the APA reference list's long
% https://doi.org/... strings are unbreakable boxes and run past the right
% margin -- the reference list is the one place this paper is full of them.
\usepackage{xurl}
\usepackage{etoolbox}

% --- Civilytics palette ---
\definecolor{paper}{HTML}{FAF7F2}
\definecolor{paper2}{HTML}{F2EDE4}
\definecolor{ink}{HTML}{0E1A2B}
\definecolor{ink2}{HTML}{2B3A52}
\definecolor{ink3}{HTML}{5A6A82}
\definecolor{rule}{HTML}{D6CEBD}
\definecolor{ruleStrong}{HTML}{B8AE97}
\definecolor{navy}{HTML}{22406A}
\definecolor{ember}{HTML}{C25311}
\definecolor{emberDark}{HTML}{923D00}
\definecolor{teal}{HTML}{1F6F70}
\definecolor{plum}{HTML}{6B3A5E}

% --- Page color ---
\pagecolor{paper}
\color{ink}

% --- Fonts (require local install or fontspec lookup) ---
% The weight suffix is "Semibold" (lowercase b) -- that is how the Source Serif 4
% family is registered with fontconfig. The original "SemiBold" fails with
% "the font Source Serif 4 SemiBold cannot be found". Note Inter spells the same
% weight "SemiBold" with a capital B, so the two are genuinely inconsistent and
% neither spelling works for both families.
\setmainfont{Source Serif 4}[
  UprightFont = *,
  ItalicFont = * Italic,
  BoldFont = * Semibold,
  BoldItalicFont = * Semibold Italic,
  Ligatures = TeX,
]
\setsansfont{Inter}[Ligatures = TeX]
\setmonofont{JetBrains Mono}[Scale = 0.92]
\newfontfamily\displayfont{Libre Franklin}[Ligatures = TeX]

% --- Hyperlinks ---
\hypersetup{
  colorlinks = true,
  linkcolor  = navy,
  urlcolor   = navy,
  citecolor  = navy,
  filecolor  = navy,
}

% --- Section headings ---
\titleformat{\section}[block]
  {\bfseries\Large\displayfont\color{ink}}
  {}{0pt}
  {\titlerule[0.4pt]\vspace{4pt}}[\vspace{-6pt}]
\titleformat{\subsection}[block]
  {\bfseries\large\displayfont\color{ink}}{}{0pt}{}
\titleformat{\subsubsection}[block]
  {\bfseries\normalsize\displayfont\color{ink}}{}{0pt}{}

% --- Page header / footer ---
\pagestyle{fancy}
\fancyhf{}
\renewcommand{\headrulewidth}{0pt}
\renewcommand{\footrulewidth}{0pt}
% Running head: set per document with \renewcommand{\cvrunninghead}{...}.
%
% Deliberately NOT \thetitle. Two reasons. First, it does not survive: \thetitle
% expands correctly at \begin{document} but renders as "0.0" in the shipped
% header, so something between there and shipout clears the title macros
% (titling documents that \maketitle discards them). Second, and more
% importantly, a full academic title is far too long for a running header --
% this paper's wraps to two lines and collides with the left-hand mark. A short
% form is the correct typography regardless.
\providecommand{\cvrunninghead}{}
\fancyhead[L]{\sffamily\scriptsize\color{ink3}\MakeUppercase{Civilytics Consulting}}
\fancyhead[R]{\sffamily\scriptsize\color{ink3}\cvrunninghead}
\fancyfoot[L]{\sffamily\scriptsize\color{ink3}civilytics.consulting}
\fancyfoot[C]{\sffamily\scriptsize\color{ink3}\thepage}
\fancyfoot[R]{\sffamily\scriptsize\color{ink3}\textcopyright\ 2026}

% --- Captions ---
\captionsetup{
  font={sf,small,color=ink3},
  labelfont={sf,bf,color=ink3},
  labelsep=period,
  justification=raggedright,
  singlelinecheck=false,
}

% --- Pullquote ---
\newtcolorbox{pullquote}{
  enhanced, breakable, frame hidden, colback=paper,
  borderline west = {2pt}{0pt}{ember},
  left=14pt, right=4pt, top=4pt, bottom=4pt,
  fontupper={\itshape\rmfamily\large\color{ink}},
}

% --- Code blocks (Pandoc + listings or minted both inherit colors below) ---
\definecolor{codebg}{HTML}{0E1A2B}
\definecolor{codefg}{HTML}{FAF7F2}

% --- Lists tighter ---
\setlist{itemsep=2pt, topsep=4pt}

% --- Tabular numerals ---
\newcommand{\tnum}[1]{{\addfontfeature{Numbers={Tabular,Lining}}#1}}

% --- Ember accent rule above section starts ---
\newcommand{\emberrule}{{\color{ember}\rule{60pt}{3pt}}\par\vspace{4pt}}

% --- APA hanging indent for the reference list ---------------------------
% The paper's references are hand-maintained markdown paragraphs, not a
% citeproc-generated bibliography, so nothing gives them a hanging indent for
% free. Wrap them with \cvrefsbegin ... \cvrefsend (raw LaTeX blocks in the
% .qmd, which non-LaTeX writers drop). Negative \parindent against a positive
% \leftskip is the standard hanging-indent idiom: first line flush left,
% continuation lines indented.
%
% NB: these are deliberately NOT named \beginrefs/\endrefs. LaTeX's
% \@ifdefinable refuses any command whose name starts with "end" -- it reserves
% that prefix for environment machinery -- and reports the refusal as the
% thoroughly misleading "Command \endrefs already defined."
\newcommand{\cvrefsbegin}{%
  \begingroup
  \setlength{\parindent}{-0.5in}%
  \setlength{\leftskip}{0.5in}%
  \setlength{\parskip}{6pt}%
  % Ragged right, not justified. Entries are short measures full of unbreakable
  % DOIs; justifying them opened visible rivers of interword space. Ragged right
  % is also the conventional setting for a bibliography.
  \raggedright
  \sloppy
}
\newcommand{\cvrefsend}{\endgroup}
